% ===== PRÉAMBULE CANONIQUE — à coller VERBATIM en tête de CHAQUE .tex =====
\documentclass[11pt,a4paper]{article}
\usepackage[utf8]{inputenc}\usepackage[T1]{fontenc}\usepackage{lmodern}
\usepackage[french]{babel}
\usepackage{amsmath,amssymb,mathtools}\usepackage{icomma}
\usepackage[version=4]{mhchem}          % \ce{...} : équations & formules
\usepackage{siunitx}\sisetup{locale=FR} % \qty{}{}, \si{}, \ang{}
\usepackage{chemfig}                    % \chemfig{...} : structures développées
\usepackage{xcolor}\usepackage{colortbl}
\usepackage{tikz}\usepackage{pgfplots}\pgfplotsset{compat=1.18}
\usepackage[most]{tcolorbox}\usepackage{enumitem}\usepackage{array,booktabs}
\usepackage[a4paper,margin=2.2cm]{geometry}
\usetikzlibrary{arrows.meta,calc,angles,quotes,positioning,patterns,backgrounds,shapes.geometric}
\definecolor{primary}{RGB}{222,49,99}\definecolor{primarydark}{RGB}{150,22,55}
\definecolor{defcol}{RGB}{0,114,178}\definecolor{methcol}{RGB}{52,92,148}
\definecolor{excol}{RGB}{0,135,90}\definecolor{attncol}{RGB}{200,40,55}
\tcbset{boxbase/.style={enhanced,breakable,boxrule=0.9pt,arc=2.5pt,
  left=3mm,right=3mm,top=2.3mm,bottom=2.3mm,fonttitle=\bfseries,coltitle=white}}
% --- boîtes TUTORIEL (noms EXACTS, ne pas renommer) ---
\newtcolorbox{cle}[1][]{boxbase,colback=defcol!6,colframe=defcol,
  title={\textsc{À retenir}\ifx\relax#1\relax\relax\else\textmd{ : \itshape #1}\fi}}
\newtcolorbox{methode}[1][]{boxbase,colback=methcol!7,colframe=methcol,sharp corners,
  title={\textsc{Méthode}\ifx\relax#1\relax\relax\else\textmd{ --- \itshape #1}\fi}}
\newtcolorbox{exemplebox}[1][]{boxbase,colback=excol!5,colframe=excol,
  title={\textsc{Exemple}\ifx\relax#1\relax\relax\else\textmd{ : \itshape #1}\fi}}
\newtcolorbox{piege}[1][]{boxbase,colback=attncol!6,colframe=attncol,
  title={\textsc{Piège}\ifx\relax#1\relax\relax\else\textmd{ : \itshape #1}\fi}}
\newtcolorbox{remarque}[1][]{boxbase,colback=black!4,colframe=black!45,
  title={\textsc{Remarque}\ifx\relax#1\relax\relax\else\textmd{ : \itshape #1}\fi}}
% --- boîtes EXERCICE / CORRIGÉ (noms EXACTS, auto-comptées) ---
\newtcolorbox[auto counter]{exo}[1][]{boxbase,colback=primary!4,colframe=primary,
  title={\textsc{Exercice}~\thetcbcounter\ifx\relax#1\relax\relax\else\textmd{ --- \itshape #1}\fi}}
\newtcolorbox[auto counter]{corrige}[1][]{boxbase,colback=excol!4,colframe=excol,
  title={\textsc{Corrigé}~\thetcbcounter\ifx\relax#1\relax\relax\else\textmd{ --- \itshape #1}\fi}}
\newcommand{\R}{\mathbb{R}}\newcommand{\N}{\mathbb{N}}\newcommand{\Z}{\mathbb{Z}}\newcommand{\ds}{\displaystyle}
% (un \usepackage{fancyhdr}+en-tête est possible ; il est ignoré côté web)
% =========================================================================
\begin{document}
% THEME_TITLE: Nombre d'oxydation
\begin{center}\color{primarydark}\Large\bfseries Thème 1 — Le nombre d'oxydation\end{center}

Une réaction d'oxydo-réduction est un \textbf{transfert d'électrons} : une espèce en cède, une
autre en capte. Pour repérer qui perd et qui gagne des électrons, on utilise le \textbf{nombre
d'oxydation} (n.o.).

\begin{cle}[ce qu'est le nombre d'oxydation]
Le \textbf{nombre d'oxydation} d'un atome est la \textbf{charge fictive} qu'il porterait si chaque
doublet de liaison était attribué \emph{entièrement} à l'atome le plus \textbf{électronégatif}. On le
note en \textbf{chiffres romains} ($+\mathrm{I}$, $-\mathrm{II}$, $+\mathrm{VI}$\ldots). C'est un
nombre \emph{formel} : il ne coïncide avec la charge réelle que pour un \textbf{ion monoatomique}.
\end{cle}

\begin{cle}[les règles d'attribution, à appliquer dans l'ordre]
\begin{enumerate}[label=\textbf{\arabic*.},leftmargin=2.0em,topsep=1pt,itemsep=1pt]
  \item \textbf{Corps simple} (\ce{O2}, \ce{H2}, \ce{Fe}, \ce{Cl2}, \ce{S}\ldots) : $\mathrm{n.o.}=0$.
  \item \textbf{Ion monoatomique} : $\mathrm{n.o.}=$ sa charge
        ($\ce{Ca^{2+}}\!\to +\mathrm{II}$, $\ce{Cl^-}\!\to -\mathrm{I}$).
  \item \textbf{Hydrogène} $=+\mathrm{I}$, \emph{sauf} dans les hydrures métalliques
        (\ce{NaH}, \ce{CaH2}) où il vaut $-\mathrm{I}$.
  \item \textbf{Oxygène} $=-\mathrm{II}$, \emph{sauf} dans les peroxydes (\ce{H2O2}, \ce{Na2O2})
        où il vaut $-\mathrm{I}$.
  \item \textbf{Halogènes} (groupe~17) $=-\mathrm{I}$ (sauf combinés à O ou entre eux) ;
        \textbf{alcalins} (gr.~1) $=+\mathrm{I}$ ; \textbf{alcalino-terreux} (gr.~2) $=+\mathrm{II}$.
\end{enumerate}
\end{cle}

\begin{cle}[la règle-clé de calcul]
La \textbf{somme} des nombres d'oxydation de tous les atomes d'une espèce est égale à sa
\textbf{charge globale} :
\[
\sum_{\text{atomes}}\mathrm{n.o.}=\text{charge globale}\qquad(\;=0\text{ pour une molécule neutre}).
\]
\end{cle}

\begin{methode}[calculer le n.o. d'un atome $X$]
\begin{enumerate}[label=\textbf{\arabic*.},leftmargin=2.0em,topsep=1pt,itemsep=1pt]
  \item Noter la \textbf{charge globale} de l'espèce ($0$ si neutre).
  \item Placer les n.o. « évidents » (O $=-\mathrm{II}$, H $=+\mathrm{I}$, halogène $=-\mathrm{I}$\ldots).
  \item Poser $x=\mathrm{n.o.}(X)$ et écrire : $\sum \mathrm{n.o.}=\text{charge globale}$.
  \item \textbf{Résoudre} l'équation en $x$.
\end{enumerate}
\end{methode}

\begin{exemplebox}[le chrome dans le dichromate \ce{K2Cr2O7}]
Molécule neutre. On a $\mathrm{n.o.}(\ce{K})=+\mathrm{I}$, $\mathrm{n.o.}(\ce{O})=-\mathrm{II}$.
Soit $x=\mathrm{n.o.}(\ce{Cr})$, avec 2 K, 2 Cr, 7 O :
\[
\underbrace{2(+1)}_{\text{K}}+\underbrace{2x}_{\text{Cr}}+\underbrace{7(-2)}_{\text{O}}=0
\;\Longrightarrow\;2x-12=0\;\Longrightarrow\;x=+6.
\]
Donc $\mathrm{n.o.}(\ce{Cr})=+\mathrm{VI}$.
\end{exemplebox}

\begin{cle}[oxydation, réduction, oxydant, réducteur]
Au cours d'une réaction, pour un atome donné :
\begin{itemize}[leftmargin=1.5em,topsep=1pt,itemsep=1pt]
  \item son n.o. \textbf{augmente} $\Rightarrow$ il est \textbf{oxydé} : il \textbf{perd} des
        électrons. L'espèce qui le contient est le \textbf{réducteur}.
  \item son n.o. \textbf{diminue} $\Rightarrow$ il est \textbf{réduit} : il \textbf{capte} des
        électrons. L'espèce qui le contient est l'\textbf{oxydant}.
\end{itemize}
Les deux phénomènes sont toujours \textbf{simultanés}.
\end{cle}

\begin{piege}[les trois erreurs classiques]
\begin{itemize}[leftmargin=1.5em,topsep=1pt,itemsep=1pt]
  \item H n'est \textbf{pas toujours} $+\mathrm{I}$ ($-\mathrm{I}$ dans les hydrures, $0$ dans \ce{H2}) ;
        O n'est \textbf{pas toujours} $-\mathrm{II}$ ($-\mathrm{I}$ dans les peroxydes).
  \item \textbf{L'oxydant est l'espèce qui est \emph{réduite}} ; le réducteur est celui qui est
        \emph{oxydé}. (Retenir : « on \emph{réduit} un oxydant, on \emph{oxyde} un réducteur ».)
  \item Un élément qui s'oxyde \textbf{perd} des électrons (n.o.\ $\nearrow$) ; ne pas confondre avec
        « capter ».
\end{itemize}
\end{piege}

\end{document}
